\documentclass[12pt,a4paper]{article}
\usepackage{amsmath}
\usepackage[margin=1in]{geometry}
\usepackage{hyperref}
\hypersetup{hidelinks}
\begin{document}

\begin{flushright}
Shuaidong Gao\\
Chongqing Institute of Foreign Studies\\
Qijiang Campus, Chongqing 401420, China\\
gsd3247186514@gmail.com\\
ORCID: 0009-0004-5641-3581\\
\today
\end{flushright}

\vspace{1cm}

\noindent\textbf{The Editor}\\
\textit{Nature}\\
4 Crinan Street\\
London N1 9XW, UK

\vspace{1cm}

\noindent Dear Editor,

\vspace{0.3cm}

I submit for your consideration my manuscript entitled \textbf{``Sample-Size Phase Transition Ends the Hard-versus-Soft Clustering Debate.''}

I am a first-year undergraduate student at Chongqing Institute of Foreign Studies, majoring in Digital Economy. This research was completed independently over the past six months, using my personal time and computational resources. As a new entrant to academia, I have no prior institutional pressures or publication expectations---this work is driven purely by curiosity.

\bigskip

\noindent\textbf{What I discovered.} When applying causal discovery to clustered data, a two-stage method (hard pre-clustering followed by causal inference) and an end-to-end method (joint clustering and causal optimization) exhibit a sharp sample-size phase transition. Hard clustering dominates at small $n$; soft clustering dominates at large $n$. The transition is governed by an empirical scaling law $n_{\text{crit}} \approx 6.3 \times d^{0.90}$ ($R^2 = 0.985$), and resolves two decades of controversy in the clustering literature.

\noindent\textbf{Why this matters beyond my field.} I replicated the identical phase transition on \textbf{six dataset types (four positive modalities and two negative controls)}:
\begin{itemize}
    \item \textbf{Cancer transcriptomics} (33 TCGA types, $r = 0.827$, $p = 3.0 \times 10^{-9}$)
    \item \textbf{Handwritten digit images} (MNIST, exponential fit $R^2 = 0.977$)
    \item \textbf{Single-cell RNA-seq} (PBMC 10k, the steepest transition observed)
    \item \textbf{Text documents} (20 Newsgroups, uniquely non-zero baseline)
    \item \textbf{Synthetic null data} (SNR $\ll$ 1) --- flat $\Delta \approx 0$ across all $n$, confirming the phenomenon requires genuine cluster structure
    \item \textbf{CIFAR-10 images} --- additional negative control, no cluster structure $\Rightarrow$ no phase transition
\end{itemize}
The monotonic decrease of gate gain $\Delta$ with increasing $n$ holds for all modalities with genuine cluster structure. The quantitative parameters---sharpness, asymptotic value, and baseline---vary by data type, but the \emph{direction} of the transition is invariant. To my knowledge, this is the first demonstration of a universal scaling law in causal machine learning that bridges genomics, computer vision, single-cell biology, natural language processing, and epigenomics.

\noindent\textbf{Why this resolves a long-standing controversy.} The hard-versus-soft clustering debate has persisted for two decades, with published benchmarks returning contradictory conclusions. I show that all 13 published benchmark datasets fall in the large-$n$ ($n \gg n_{\text{crit}}$) regime---the soft-clustering advantage zone---systematically omitting the small-$n$ regime where hard clustering dominates and most biomedical applications operate. The ``controversy'' is a disagreement about which sample-size regime different studies happened to examine.

\noindent\textbf{Why \textit{Nature}.} This work:
\begin{enumerate}
    \item Reports a novel, quantitatively confirmed scientific law ($n_{\text{crit}}$ formula, $R^2 = 0.985$)
    \item Validates it across six dataset types from completely different disciplines, including two negative controls
    \item Resolves a twenty-year controversy in the machine learning literature
    \item Provides a practical decision rule ($n/d \approx 4$) with immediate applications in biomedical research
\end{enumerate}

\noindent\textbf{Data and code availability.} All data are publicly available (TCGA, 10x Genomics, MNIST, 20 Newsgroups). CIFAR-10 was used as a negative control only. Full experimental results---including the complete 33-cancer $d{=}200$ 10-seed sweep, within-cancer resampling, single-cell, MNIST, text, and methylation benchmarks---and a complete replication package are available for peer review upon request, and will be made publicly available on GitHub upon acceptance and completion of patent filing.

\noindent\textbf{Competing interests.} A Chinese patent application covering the Cluster-Aware Gating mechanism has been filed. I declare no other competing interests.

\vspace{1cm}
\noindent I look forward to your consideration.

\vspace{0.5cm}
\noindent Sincerely,\\
\vspace{0.3cm}
\noindent Shuaidong Gao\\
\noindent Chongqing Institute of Foreign Studies

\end{document}
